\documentclass[aspectratio=169,12pt]{beamer}
\usepackage{fontspec}
\usepackage[portuguese]{babel}
\usepackage{xcolor,listings,tikz,booktabs,amsmath,amssymb,graphicx,multicol,array,colortbl}
\usetikzlibrary{arrows.meta,positioning,shapes.geometric,calc}

\definecolor{navy}{HTML}{0D1B3E}
\definecolor{gold}{HTML}{F5A623}
\definecolor{qgreen}{HTML}{1E8B4C}
\definecolor{qblue}{HTML}{1A6EAE}
\definecolor{codebg}{HTML}{EEF3F8}

\usetheme{default}
\useinnertheme{circles}
\setbeamertemplate{navigation symbols}{}
\setbeamertemplate{footline}{%
  \hfill{\color{gray!60}\scriptsize\insertframenumber/\inserttotalframenumber}\hspace{.6em}\vspace{.4em}}

\setbeamercolor{frametitle}{bg=navy,fg=white}
\setbeamercolor{structure}{fg=navy}
\setbeamercolor{alerted text}{fg=gold}
\setbeamercolor{item}{fg=gold}
\setbeamercolor{subitem}{fg=navy}
\setbeamercolor{block title}{bg=qblue,fg=white}
\setbeamercolor{block body}{bg=qblue!8,fg=black}
\setbeamercolor{block title alerted}{bg=gold,fg=navy}
\setbeamercolor{block body alerted}{bg=gold!15,fg=navy}
\setbeamercolor{block title example}{bg=qgreen,fg=white}
\setbeamercolor{block body example}{bg=qgreen!10,fg=black}
\setbeamerfont{frametitle}{size=\large,series=\bfseries}
\setbeamerfont{block title}{series=\bfseries}

\setbeamertemplate{title page}{%
  \begin{tikzpicture}[remember picture,overlay]
    \fill[navy](current page.north west)rectangle(current page.south east);
    \fill[gold]([yshift=.7cm]current page.south west)rectangle(current page.south east);
    \node[anchor=south,yshift=.73cm,navy,font=\tiny\bfseries]at(current page.south)
      {INTENSIVÃO QUANT AI \textbullet{} IMPACT UFSCAR};
  \end{tikzpicture}
  \vspace{.65cm}
  \begin{center}
    {\color{gold}\scriptsize\bfseries INTENSIVÃO QUANT AI \textbullet{} IMPACT UFSCAR}\\[.4cm]
    {\color{white}\LARGE\bfseries\inserttitle}\\[.25cm]
    {\color{gold!85}\normalsize\insertsubtitle}\\[.5cm]
    {\color{white!70}\small\insertauthor}\\[.1cm]
    {\color{white!50}\footnotesize\insertdate}
  \end{center}}

\lstdefinestyle{python}{language=Python,
  basicstyle=\ttfamily\scriptsize\color{qblue},
  keywordstyle=\color{navy}\bfseries,
  stringstyle=\color{qgreen!80!black},
  commentstyle=\color{gray!70}\itshape,
  backgroundcolor=\color{codebg},
  frame=l,framerule=2pt,rulecolor=\color{gold},
  xleftmargin=1em,breaklines=true,showstringspaces=false,tabsize=4,
  extendedchars=false,inputencoding=ascii}
\lstset{style=python}

\author{Felipe Maldonado\\{\scriptsize VP Diretoria Quant~\textbullet{}~Impact UFSCAR}}
\date{Intensivão Quant AI 2026}
\title{Aula 05 --- Backtest v1}
\subtitle{Construindo e Avaliando o Primeiro Backtest}

\begin{document}

\begin{frame}[plain]
\titlepage
\end{frame}

\begin{frame}[fragile]{Nesta Aula}
\begin{columns}[T]
  \column{0.5\textwidth}
\begin{block}{Teoria (20 min)}
\begin{itemize}
  \item Hipóteses do backtest: o que estamos testando?
  \item Métricas de performance: Sharpe, Sortino, MDD, Calmar
  \item O problema do look-ahead bias: shift(1)
  \item A importância de usar pesos realistas
\end{itemize}
\end{block}
  \column{0.47\textwidth}
\begin{block}{Código (40 min)}
\begin{itemize}
  \item \texttt{construir\_portfolio()} — top 10\%, weights.shift(1)
  \item \texttt{calcular\_metricas()} — todas as métricas de uma vez
  \item Equity curve e drawdown chart
  \item Análise de retornos por ano (heatmap)
  \item Salvar \texttt{pesos\_v1.parquet} e \texttt{retorno\_carteira.parquet}
\end{itemize}
\end{block}
\end{columns}
\end{frame}

\begin{frame}[fragile]{Métricas de Performance}
\begin{columns}[T]
  \column{0.5\textwidth}
\textbf{Sharpe Ratio:}
\[ S = \frac{\overline{r} - r_f}{\sigma_r} \times \sqrt{12} \]
\vspace{.1em}
\textbf{Sortino Ratio:}
\[ \text{Sort} = \frac{\overline{r} - r_f}{\sigma_{\text{neg}}} \times \sqrt{12} \]
$\sigma_{\text{neg}}$: desvio padrão apenas dos retornos negativos.
\vspace{.2em}
\textbf{Maximum Drawdown (MDD):}
\[ \text{MDD} = \min_t \frac{V_t - \max_{s \leq t} V_s}{\max_{s \leq t} V_s} \]
\textbf{Calmar Ratio:}
\[ \text{Calmar} = \frac{\text{CAGR}}{|\text{MDD}|} \]
  \column{0.47\textwidth}
\begin{block}{Referências de Mercado}
\begin{itemize}
  \item Sharpe $> 1{,}0$: excelente (difícil na realidade brasileira)
  \item Sharpe $> 0{,}5$: bom para estratégias long-only no IBOV
  \item MDD $< 20\%$: confortável para a maioria das gestoras
  \item Calmar $> 0{,}5$: relação risco/retorno aceitável
\end{itemize}
\end{block}\vspace{.3em}
\begin{exampleblock}{CAGR — Retorno Anualizado}

\[ \text{CAGR} = \left(\prod_{t=1}^{T}(1+r_t)\right)^{12/T} - 1 \]
Compõe retornos mensais e anualiza pelo número de meses $T$.
\end{exampleblock}
\end{columns}
\end{frame}

\begin{frame}[fragile]{Estrutura do Backtest}
\begin{columns}[T]
  \column{0.5\textwidth}
\textbf{Por que shift(1)?}
\vspace{.3em}
\begin{itemize}
  \item Sinal calculado no \textbf{final do mês $t-1$}
  \item Portfólio formado no \textbf{início do mês $t$}
  \item Retorno realizado no \textbf{mês $t$}
\end{itemize}
\vspace{.3em}
\begin{lstlisting}
# ERRADO: look-ahead bias
ret_carteira = (sinal > 0) * ret_mensais

# CORRETO: shift(1) evita o bias
pesos = construir_pesos(sinal)
ret_carteira = (pesos.shift(1) * ret_mensais).sum(axis=1)
\end{lstlisting}
  \column{0.47\textwidth}
\begin{alertblock}{Hipóteses do backtest}
\begin{itemize}
  \item Execução ao preço de fechamento do mês
  \item Liquidez infinita (sem market impact inicial)
  \item Sem short: long-only com equal-weight no top 10\%
  \item Benchmark: IBOVESPA (BOVA11)
\end{itemize}
\end{alertblock}
\end{columns}
\end{frame}

\begin{frame}[fragile]{O que Vamos Construir}
\begin{columns}[T]
  \column{0.52\textwidth}
\begin{block}{Funções a implementar}
\begin{itemize}
  \item \texttt{construir\_portfolio(sinal, ret, n\_pct=0.1)} $\to$ retornos da carteira
  \item \texttt{calcular\_metricas(ret, rf=0.1025)} $\to$ dicionário de métricas
  \item \texttt{exibir\_metricas(metricas)} $\to$ tabela formatada
  \item \texttt{plot\_equity\_drawdown(ret)} $\to$ equity curve + MDD
\end{itemize}
\end{block}
  \column{0.45\textwidth}
\begin{block}{Entradas (parquets)}
\begin{itemize}
  \item \texttt{sinal\_v1.parquet}
  \item \texttt{retornos\_mensais\_limpo.parquet}
\end{itemize}
\end{block}
\vspace{.4em}
\begin{exampleblock}{Saídas (parquets)}
\begin{itemize}
  \item \texttt{pesos\_v1.parquet}
  \item \texttt{retorno\_carteira.parquet}
\end{itemize}
\end{exampleblock}
\end{columns}
\end{frame}

\begin{frame}[fragile]{Referências}
\begin{multicols}{2}
\tiny
\begin{itemize}
  \item \textbf{Sharpe, W.} (1966). Mutual fund performance. \textit{Journal of Business}, 39(1), 119--138.
  \item \textbf{Sortino, F. \& van der Meer, R.} (1991). Downside risk. \textit{Journal of Portfolio Management}, 17(4), 27--31.
  \item \textbf{Jegadeesh, N. \& Titman, S.} (1993). Returns to buying winners. \textit{Journal of Finance}, 48(1), 65--91.
  \item \textbf{Bailey, D. \& López de Prado, M.} (2014). The deflated Sharpe ratio. \textit{JPM}, 40(5).
\end{itemize}
\end{multicols}
\end{frame}


\end{document}
